% Part: proof-theory
% Chapter: natural-deduction
命题逻辑的语言~$\Lang L_0$ 由以下符号构成：无穷多个不同的\emph{命题变元}：$p_0$, $p_1$, $p_2$, $\dots$；以及五个\emph{联结词}：$\lnot$, $\land$, $\lor$, $\lif$, $\liff$。这些符号（命题变元与联结词）都互不相同。
% Section: rules-N2
我们通常用 $p$, $q$, $r$, \dots 来表示命题变元。

\begin{explain}
$\Lang L_0$ 的任意有限符号序列都只是一个符号串。
\emph{构造序列}是一个符号串的有限序列，其中每个串都是根据某些形成规则从前面的串构造出来的。
一个\emph{合式公式}（或简称\emph{公式}）是某个构造序列中的最后一个串。
\end{explain}

公式归纳定义如下：
\begin{enumerate}
\item （归纳基）每个命题变元~$p$ 都是一个公式。
\item （归纳步）若~$!A$ 和~$!B$ 是公式，则下列也是公式：
\begin{align*}
& \lnot !A && \text{（$!A$ 的否定）}\\
& (!A \land !B) && \text{（$!A$ 与 $!B$ 的合取）}\\
& (!A \lor !B) && \text{（$!A$ 与 $!B$ 的析取）}\\
& (!A \lif !B) && \text{（$!A$ 与 $!B$ 的条件句）}\\
& (!A \liff !B) && \text{（$!A$ 与 $!B$ 的双条件句）}
\end{align*}
\end{enumerate}
除此之外，没有别的东西是公式。

注意，符号 $\lif$ 和 $\land$ 以及 $\lor$ 和 $\liff$ 都是\emph{二元中缀运算符}：它们写在两个公式之间。
为明确起见，在本书中我们将它们当作二元联结词，但在定义真值函数时，我们允许它们接收多个参数。
例如，$\lif(A_1, \dots, A_n, B)$ 只是 $A_1 \lif (A_2 \lif \dots (A_n \lif B)\dots)$ 的简写，而 $\land(A_1, \dots, A_n)$ 是 $(\dots (A_1 \land A_2) \land \dots \land A_n)$ 的简写，但 \texttt{if}

项或公式的\emph{高度}是一个递归定义的自然数。定义如下：
\begin{enumerate}
\item 对任意常项符号~$c$、变元~$x$ 以及原子公式~$\Atom{R}{t_1, \dots, t_n}$，定义 $\height{c} = \height{x} = \height{\Atom{R}{t_1, \dots, t_n}} = 1$。
\item 若 $t_1$, \dots, $t_n$ 是项且 $f$ 是一个 $n$元函数符号，则 $\height{f(t_1, \dots, t_n)} = \max(\height{t_1}, \dots, \height{t_n}) + 1$。
\item 若 $!A_1$, \dots, $!A_n$ 是公式且 $\star$ 是一个 $n$元逻辑算子，则 $\height{\star(!A_1, \dots, !A_n)} = \max(\height{!A_1}, \dots, \height{!A_n}) + 1$。
\end{enumerate}
换句话说，一个项或公式的高度，比它的各个直接真子项或直接真子公式（视情况而定）的高度的最大值大~1。如果一个项或公式没有直接真子项或直接真子公式（即它是常项、变元或原子公式），则其高度为~1。

等价地，$\height{u}$ 是满足以下条件的最大自然数~$n$：存在一个长度为~$n$ 的、$u$ 的构造序列（参见\olref[fol][syn][seq]{defn:formation-sequence}）。任何 $u$ 的构造序列的长度都至少是 $\height{u}$，并且存在某个 $u$ 的构造序列，其长度恰好是 $\height{u}$。

\documentclass[../../../include/open-logic-section]{subfiles}

\begin{document}

\begin{table}
\begin{tabular}{@{}c@{}c@{}}
\multicolumn{2}{l}{\textbf{公理：}}\\
\multicolumn{2}{c}{$x:!A \Sequent x:!A$}\\[3ex]
\multicolumn{2}{l}{\textbf{推理规则：}}\\[3ex]
\multicolumn{2}{c}{\Axiom$x: !A, \Gamma \fCenter N:!B$
\RightLabel{\Intro{\lif}}
\UnaryInf$\Gamma \fCenter \lambd[\typeof{x}{!A}][N] : !A \lif !B $
\DisplayProof}
\\[4ex]
\multicolumn{2}{c}{\Axiom$\Gamma_1 \fCenter N:!A \lif !B$
\Axiom$\Gamma_2 \fCenter M:!A$
\RightLabel{\Elim{\lif}}
\BinaryInf$\Gamma_1, \Gamma_2 \fCenter (NM) : !B$
\DisplayProof}
\\[4ex]
\multicolumn{2}{c}{\Axiom$\Gamma_1 \fCenter N_1 : !A$
\Axiom$\Gamma_2 \fCenter N_2 : !B$
\RightLabel{\Intro{\land}}
\BinaryInf$\Gamma_1, \Gamma_2 \fCenter \pair{N_1}{N_2} : !A \land !B $
\DisplayProof}
\\[4ex]
\Axiom$\Gamma \fCenter N: !A \land !B$
\RightLabel{\Elim{\land}[1]}
\UnaryInf$\Gamma \fCenter \proj{1}{N} : !A$
\DisplayProof
&
\Axiom$\Gamma \fCenter N:!A \land !B$
\RightLabel{\Elim{\land}[2]}
\UnaryInf$\Gamma \fCenter \proj{2}{N} : !B$
\DisplayProof\\[3ex]
\\[4ex]
\Axiom$\Gamma \fCenter N: !A$
\RightLabel{\Intro{\lor}[1]}
\UnaryInf$\Gamma \fCenter \inj[!B]{1}{N} : !A \lor !B$
\DisplayProof
&
\Axiom$ \Gamma \fCenter N : !B$
\RightLabel{\Intro{\lor}[2]}
\UnaryInf$\Gamma \fCenter \inj[!A]{2}{N} : !A \lor !B$
\DisplayProof\\[3ex]
\multicolumn{2}{c}{
\begin{tabular}[b]{@{}l@{}}
\Axiom$\Gamma_1 \fCenter M : !A \lor !B$
\Axiom$x:!A, \Gamma_2 \fCenter N_1 : !C$
\Axiom$y:!B, \Gamma_3 \fCenter N_2: !C$
\RightLabel{\Elim{\lor}}
\TrinaryInf$\Gamma_1, \Gamma_2, \Gamma_3 \fCenter
\dcase{M}{x}{N_1}{y}{N_2} : !C$
\DisplayProof
\end{tabular}
}
\\[4ex]
\multicolumn{2}{c}{\Axiom$\Gamma \fCenter N: \lfalse$
\RightLabel{\FalseInt}
\UnaryInf$\Gamma \fCenter \abort{!A}{N} : !A$
\DisplayProof}
\end{tabular}
\caption{命题项标记直觉主义自然演绎系统 \Log{tN2ip} 的规则。}
\ollabel{tab:tN2ip}
\end{table}

\end{document}